% Version 0.27. Actual symplectic Model construction on the original datum.
% Pinned Lean source: 7cbf7a0c32fca1badf05931879a37c3f8a2b03b4.
\section{Constructing the actual symplectic deformation model}
\label{sec:symplectic-model-construction}

\status{MX1--MX10 reconstruct every field of the source's
\texttt{Model D}.  The intrinsic constructor takes the original
\texttt{Good D}, \texttt{PhiCond D.Gam}, \texttt{D.Nums} and an
actual \texttt{LocBase D}; FP8 supplies its universal
\texttt{LemmaP}.  Substitution of the exact arithmetic suppliers
DX2, DX4 and DX8 gives a model for every original good datum.
The construction does not supply image density, automorphy,
rigidity, a finite trace catalogue or the later geometric
comparison of points.}

\subsection{The objects that must be kept fixed}

Fix an algebraic closure $\Omega$ of $\mathbb Q$, an odd prime
$\ell$, the specified coefficient DVR $O=c.O$, its fraction field
$E$ and finite residue field $k$, and an original
$D:\texttt{Raw}\,\Omega\,c$ satisfying \texttt{Good}.
No field, coefficient ring, character or decomposition group is
changed in this chapter.  Put $e=D.e\ge2$.
Write $G=\operatorname{Gal}(\Omega/\mathbb Q)$,
$G_p=D.Gp(p)$, $I_p=D.Ip(p)$ and
\[
 \Psi=\chi_\ell^{e}\psi,\qquad
 \alpha=\chi_\ell\omega^{-1}\xi,\qquad
 \vartheta=\alpha\Psi,\qquad
 r_0=\operatorname{diag}(\vartheta,\Psi,\alpha,1),
 \qquad \bar r=r_0\bmod\mathfrak m_O.
 \tag{MX.1}
\]
Here $\chi_\ell$ is the integral cyclotomic character.  Set
$S_b=D.Sb$, $T=\{\ell\}\cup S_b$, and $b=|S_b|$.
The four vector weights and relevant matrix weights are
\[
 \begin{gathered}
 w_0=(1,1),\quad w_1=(0,1),\quad
 w_2=(1,0),\quad w_3=(0,0),\qquad w_{ij}=w_i-w_j,\\
 A=(1,0),\quad B=(-1,1),\quad P=(0,1),\quad H=(1,1),\\
 \mathcal W=\{0,\pm A,\pm B,\pm P,\pm H\}.
 \end{gathered}
 \tag{MX.2}
\]
These are exactly \texttt{w4}, \texttt{wt}, \texttt{wAl},
\texttt{wBe}, \texttt{wPs}, \texttt{wTh} and \texttt{wts}.
The letter $P$ in a weight means $(0,1)$; a ring denoted
$\mathcal P$ below is different.

Use the alternating form
\[
 J=\begin{pmatrix}0&0&0&1\\0&0&1&0\\0&-1&0&0\\-1&0&0&0\end{pmatrix},
 \qquad X^*=J^{-1}X^tJ,
 \qquad \mathfrak {sp}_4(O)=\{Y:Y^*=-Y\}.
 \tag{MX.3}
\]
Thus \texttt{IsGSp}$(\nu,g)$ means $g^*g=\nu I$.
In particular the form convention is the original one, including
the signs of the two entries of
$M=E_{20}+E_{31}\in\mathfrak {sp}_4(O)$, of weight $-P$.

Write $\mathcal C_O$ for the complete Noetherian local
$O$-algebras with the specified residue field $k$.
It is exactly \texttt{IsCO}; algebra maps between these rings are
automatically local and continuous by FP1.  A lift has locally
constant entries modulo every power of the maximal ideal.
For an augmented ring $x:R\to O$, an $O$-valued point derivation
satisfies
\[
 \delta(fg)=x(f)\delta(g)+x(g)\delta(f).
 \tag{MX.4}
\]
It has weight $j$ if it vanishes on every homogeneous piece of
weight different from $j$.  This convention evaluates on the
weight-$j$ piece; it does not reverse the sign of $j$.

We use a \emph{weighted presentation} in the exact sense of FP:
a surjection $O[[z_s:s\in\Sigma]]\to R$ with finite $\Sigma$
and homogeneous kernel.  Its weight pieces $R_j$ are images of
homogeneous formal series, and $R^{\mathrm{gr}}=\sum_jR_j$ consists
of their finite sums.  This subalgebra need not be the whole
complete ring.  An augmented weighted presentation always uses
the source's constant-coefficient origin: $xW$ is evaluation
at zero, so every presentation variable has origin zero.
The augmentation of $R^{\mathrm{gr}}$ is the restriction of $x$.

The input \texttt{LocBase D} is an actual ring
$B_0\in\mathcal C_O$, with a weighted presentation, origin
$x_B:B_0\to O$, and a universal family $(\rho_v)_{v\in T}$
of the source's local equations.  At $\ell$ these are the
fixed multiplier and \texttt{Ord} projector equations; at
$p\in S_b$ they are the fixed multiplier and identity on
$I_p\cap\ker\psi$.  Universality ranges over all rings of
$\mathcal C_O$, on the original $G_p$.
Its entries have weight $w_{ij}$, their origins are
$r_0(g)_{ij}$, and its prescribed smooth-point dimensions are
\[
 n_B(j)=d_\ell(j)+b\,d_p(j),
 \tag{MX.5}
\]
where $d_\ell$ is $2$ on $0,A,B,P,H$, $1$ on the other
four nonzero weights, and $0$ elsewhere; $d_p$ is $2$ at $0$,
$1$ at every nonzero weight in $\mathcal W$, and $0$ elsewhere.
The precise meaning of this smooth-point assertion is explained
and preserved in MX4--MX5.

\subsection{The actual framing problem}

% BEGIN MX1
\begin{proposition}[MX1: diagonal base point and original quotient descent]
\label{claim:MX1}
Let $\Gamma=D.Gam$ be the quotient of $G$ by the closed normal
subgroup generated by all inertia outside $D.S$.
The original $r_0,\bar r$ descend to $\Gamma$, with the form and
weight normalizations needed by FP8.  On $\Gamma$, imposing
identity on
\[
 N_q=I_q\cap\ker\bar r\ \le G_q,
                       \qquad q\in D.S\setminus T,
 \tag{MX.6}
\]
is equivalent, after pullback to $G$, to the literal
\texttt{D.UnrT} condition.  This is an equivalence of actual
continuous lifts and preserves the prescribed multiplier.
The original diagonal local family satisfies all
\texttt{LocEqs} equations, and these equations, multiplier,
normalization and unramified conditions commute with coefficient
maps in $\mathcal C_O$.
\end{proposition}
\begin{proof}
AX6 supplies the literal \texttt{facts} consequences of
\texttt{Good} on this same datum.  In particular the four
characters in (MX.1) are adically continuous, $\ker\bar r$
is open, all four characters are trivial on inertia outside
$D.S$, and
\[
 g\in I_A,\ A\text{ away from }T,\ \bar r(g)=1
           \quad\Longrightarrow\quad r_0(g)=1.
 \tag{MX.7}
\]
Here $A$ ranges over the original valuation subrings of
$\Omega$, exactly as in \texttt{UnrT}.  There is no replacement
by a convenient finite list of completions.
The kernel of any continuous lift into a complete ring is
closed: it is the intersection of the clopen conditions
$r(g)\equiv1\pmod {\mathfrak m^n}$.  Thus $r_0,\bar r$
factor through the indicated closed normal quotient.
The quotient is profinite.  For example its points are separated
by the finite quotients of $G$ containing the closed normal
subgroup; compactness then gives its inverse-limit description.
Factorization preserves adic continuity because the quotient
map is a quotient map on each finite matrix-valued reduction.

We check the stronger correspondence for arbitrary lifts.
If a lift on $G$ satisfies \texttt{UnrT}, then it kills inertia
outside $D.S$: such inertia is away from $T\subset D.S$,
and its residual image is one by AX6.  Its closed kernel
therefore contains the defining subgroup of $\Gamma$.
Conversely, a lift on $\Gamma$ satisfying (MX.6) kills the
required residual-trivial inertia at every original valuation.
Outside $D.S$ this follows from the quotient definition.
Over $q\in D.S\setminus T$, conjugacy of extensions of the
$q$-adic valuation identifies its inertia with a conjugate of
the specified $I_q$; both the representation kernel and
$\ker\bar r$ are normal.  At the trivial valuation the inertia
is trivial.  A valuation cannot lie over two different rational
primes, by B\'ezout's identity.  These cases are exhaustive,
and prove the asserted equivalence and its inverse.

The diagonal identities
$\vartheta\cdot1=\Psi\alpha$ give
$r_0(g)^*r_0(g)=\vartheta(g)I$.
The matrices are diagonal, so the block and equal-weight clauses
of \texttt{Normalises} hold over both $O$ and $k$.
For any lift $r$ of $\bar r$, the condition
$r(g)r_0(g)^{-1}\in\operatorname{Sp}_4$ with reduction one
is precisely the same fixed multiplier condition.  This proves
the form comparison used in descent.

The origin of the given local base sends its actual equations
to the diagonal local tuple.  Thus $r_0|_{G_v}$ satisfies
\texttt{LocEqs} without an additional existence assumption.
All its defining equations are polynomial: characteristic
polynomials and the projectors
$\prod_{j\ne i}(X-\psi_j(s_0)I)$, their products in
\texttt{Ord}, the multiplier equations, and prescribed
matrix equalities on inertia.  Polynomial evaluation commutes
with every algebra map.  So do the two \texttt{hEl}
equations and the block equations in \texttt{Nrm}.
Every equality in \texttt{UnrT} is preserved as well.
In particular an \texttt{OnComp} witness maps to the stated
local equation; the domain witness is only needed when the
stronger component condition itself is asserted.
\end{proof}
% END MX1

% BEGIN MX2
\begin{proposition}[MX2: the entire kill-inertia local condition]
\label{claim:MX2}
For every $q\in D.S\setminus T$, the condition that a local
lift kill $N_q$ satisfies all five \texttt{IsLocCond} axioms
used in FP8, on the original group $G_q$.  The original
diagonal local representation satisfies this condition.
\end{proposition}
\begin{proof}
Set $H=G_q$, $I=I_q$ and $N=I\cap\ker\bar r$.
These are closed normal subgroups where appropriate.
The finite group $I/N$ embeds by $\bar r$ into the diagonal
group $(k^\times)^4$, so its order is prime to $\ell$.
The original decomposition group satisfies
$H/I\simeq\widehat{\mathbb Z}$, with arithmetic Frobenius
mapping to $1$.  This is the finite unramified-extension
theorem, transported from the completion as in KB6; it concerns
the specified $D.pl(q)$, not a new local place.\footnote{J. S. Milne,
\emph{Algebraic Number Theory}, version 3.08, Proposition 7.50,
Corollary 7.52 and Theorem 7.58(a), with the finite-residue-field
discussion following Remark 7.53;
\url{https://www.jmilne.org/math/CourseNotes/ANT.pdf}.
Taking compatible finite unramified extensions gives the stated
inverse-limit group.  KB6 supplies the comparison with the
original decomposition group.}

For every finite continuous $k[H/N]$-module $V$ we have
\[
                    H^2_{\mathrm{cts}}(H/N,V)=0.
 \tag{MX.8}
\]
Here is the needed specialization in detail.  Averaging over
the finite normal subgroup $I/N$ is possible because its order
is invertible in $k$, and makes its positive-degree cohomology
zero.  The Hochschild--Serre cochain filtration therefore
identifies the group in (MX.8) with
$H^2_{\mathrm{cts}}(\widehat{\mathbb Z},V^{I/N})$.
The latter vanishes by a direct finite calculation.  A
continuous cocycle and its action factor through a finite
cyclic quotient $C_m$.  A normalized extension representing
its class is determined by a lift $u$ of the generator and
the element $u^m$ in the invariant coefficient group.
Pulling back to $C_{\ell m}$ replaces this element by
$u^{\ell m}=\ell u^m=0$.  The resulting extension splits,
so the class dies in the direct limit.  Compactness ensures
that all the cochains here factor through finite quotients;
the finite-group cochain filtration and its averaging
contraction pass to that filtered limit.  This proves (MX.8)
for exactly the locally constant cochains that the source uses.
The general filtration theorem is the classical theorem of
Hochschild--Serre; only its finite-normal-subgroup specialization
is used here.\footnote{G. Hochschild and J.-P. Serre,
\emph{Cohomology of group extensions}, Trans. Amer. Math. Soc.
74 (1953), 110--134, Chapter II, Theorems 1--2
(the finite cochain filtration and its second page);
\url{https://www.sas.rochester.edu/mth/sites/doug-ravenel/otherpapers/hochschild-serre.pdf}.
The passage to continuous finite cochains and the cyclic
vanishing calculation are given above.}

Four local-condition axioms are immediate with their full
quantifiers.  Closedness is given by the ideal generated by
the entries of $r(n)-I$, $n\in N$.  This ideal has precisely
the required universal base-change property.  Conjugation
and diagonal torus conjugation preserve identity on $N$.
If one of two lifts is one on $N$, the other is one there
exactly when their restrictions to $N$ agree; this is the
literal tangent axiom, including its reverse implication.
The diagonal base tuple kills $N$ by (MX.7).

For the remaining lifting axiom take an arbitrary surjection
$p:A\to A'$ in $\mathcal C_O$ with kernel $J_0$ satisfying
$\mathfrak m_AJ_0=0$.  Thus $J_0^2=0$ and $J_0$ is a finite
$k$-vector space.  Let $r'$ be a prescribed local lift of
the fixed multiplier which kills $N$.
Apply the form-preserving matrix lift construction of FP4 to
$r'(g)r_0(g)^{-1}$, as a function on $H/N$.
It gives continuous matrices $\widetilde r(g)$ of the same
multiplier, reducing to $r'(g)$ and constant on $N$-cosets;
choose $\widetilde r(1)=I$.
For clarity, continuity requires no Artinian hypothesis on
$A$.  The finite kernel $J_0$ has zero intersection with
$\mathfrak m_A^n$ for some $n$; matrix lifts can be chosen
consistently on the corresponding clopen residue classes,
as in FP4.  The first-order correction of a form error uses
$1/2\in O$ and preserves the fixed multiplier.

The multiplicative error
\[
 \kappa(g,h)=\widetilde r(g)\widetilde r(h)
                      \widetilde r(gh)^{-1}-I
       \in\mathfrak {sp}_4(k)\otimes_kJ_0
 \tag{MX.9}
\]
is a locally constant two-cocycle for the residual adjoint
action.  Associativity gives its cocycle identity, because
all products of two kernel entries vanish.  That action is
trivial on $N$ and continuous.  Equation (MX.8) gives a
locally constant $a:H/N\to\mathfrak {sp}_4(k)\otimes J_0$
with $da=\kappa$.  Since $\kappa$ is normalized, $a(1)=0$.
Then
\[
                   r(g)=(I-a(g))\widetilde r(g)
 \tag{MX.10}
\]
is multiplicative: its error is $\kappa-da$.
It is continuous, has the same multiplier, reduces to $r'$,
and is identity on $N$.  The linear form condition on $a$
ensures preservation of the form to first order, and
$J_0^2=0$ makes that preservation exact.  This is the
required lift for every small surjection and proves all five
axioms on the actual local group.
\end{proof}
% END MX2

% BEGIN MX3
\begin{proposition}[MX3: the actual universal framed ring and its origin]
\label{claim:MX3}
Assume \texttt{PhiCond D.Gam} and the literal \texttt{D.Nums},
and fix the actual local base $B_0$ above.
There is a full source record \texttt{Fr D L}.
In particular there are $U\in\mathcal C_O$, $i:B_0\to U$,
a continuous global lift $r_U$ and matrices $\alpha_v$
which are universal for
\[
 \begin{gathered}
 r_U(g)^*r_U(g)=\vartheta(g)I,\qquad
 \texttt{UnrT}(r_U,1),\qquad
 \alpha_v^*\alpha_v=I,\quad \bar\alpha_v=I,
 \quad\alpha_\ell=I,\\
 r_U(g)\alpha_v=\alpha_v i(\rho_v(g))
                                      \quad(g\in G_v).
 \end{gathered}
 \tag{MX.11}
\]
Let $a_j$ be the number of new relative variables of weight
$j$.  The relative presentation has exactly
\[
 a_j+d_j\quad\hbox{relations of weight }j,
 \qquad
 d_0=2,\quad d_P=d_{-P}=1,\quad d_j=0\text{ otherwise},
 \tag{MX.12}
\]
and $a_j=0$ for $j\notin\{0,P,-P\}$.
Its origin sends $r_U$ to the original $r_0$, all
$\alpha_v$ to $I$, and $i$ to $x_B$.
\end{proposition}
\begin{proof}
Apply FP8 with the group $\Gamma$ and the descended base
representations of MX1, the four weights in (MX.2), the form
(MX.3), the original local maps $G_v\to\Gamma$, distinguished
place $\ell$, and the kill-$N_q$ conditions of MX2.
The one block and its trivial involution give exactly the
source's \texttt{form O}; $w_i+w_{3-i}=(1,1)$ wherever
$J_{i,3-i}\ne0$.  MX1 checks both normalizers and all
\texttt{InH}/multiplier comparisons.  MX2 supplies every
local-condition axiom and its value at the base tuple.
The given local base supplies the original $B_0$, the
homogeneous entries, their local equations and their
fixed multiplier.  Its presentation variables have origin
zero, hence belong to $\mathfrak m_{B_0}$.
All hypotheses of FP8 have now been installed explicitly.

The obstruction finiteness required there is the first
clause of \texttt{Nums}, applied to the descended residual
representation and its specified normalizer.  FP8 makes the
number of new variables of weight $j$ equal to the
dimension of the source's unquotiented framed space
$H^1_F(\operatorname{ad}_j)$, and the number of relations
equal to $\dim_k\operatorname{Ob}_F(\operatorname{ad}_j)$.
The other two \texttt{Nums} clauses give exactly (MX.12)
and the stated vanishing of $a_j$.  These are equalities,
not estimates with a possibly larger obstruction space.
Each new variable is an actual entry of $r_U-r_0$ or
$\alpha_v-I$, with its original weight $w_{ij}$.

Universality of FP8 descends and pulls back along the
surjection $G\to\Gamma$ by MX1.  Equalities of pulled-back
representations are equivalent to equalities on $\Gamma$;
the local maps and all frame matrices are unchanged.
This is exactly universality in (MX.11) for every target
ring in $\mathcal C_O$, and not only for finite targets.

The tuple $(x_B,r_0,(I)_v)$ satisfies (MX.11), so universality
gives $x_U:U\to O$ with those values.  It kills every old
presentation variable because $x_B$ does, and every new one
because it is a centered entry of this tuple.  The composite
with the formal-series presentation is therefore constant
coefficient, by continuity and finite Taylor approximation.
Consequently its kernel has zero constant term and
$x_U$ is the presentation's actual origin.
All asserted origin identities follow from that one
universal map.  This proves every field of \texttt{Fr}.
\end{proof}
% END MX3

\subsection{Charts, independent cuts and a genuine domain}

% BEGIN MX4
\begin{proposition}[MX4: finite jets, relative smooth covers and derivation lifts]
\label{claim:MX4}
Let $x:R\to O$ be the constant-coefficient origin of a finite
weighted presentation, and put
$I=\ker x$.  Then, with $I_{\mathrm{gr}}=I\cap R^{\mathrm{gr}}$,
\[
 R^{\mathrm{gr}}/I_{\mathrm{gr}}^n\ \simeq\ R/I^n
                         \qquad(n\ge1).
 \tag{MX.13}
\]
Charts and smooth-point dimensions therefore pass between
the homogeneous subalgebra and the complete ring.
For the actual framed relative presentation of MX3 there
is a cover
\[
 B_0\ \longrightarrow\ \mathcal P=B_0[[u_t:t\in\tau]]
             \mathop{\longrightarrow}^{p}\ U
 \tag{MX.14}
\]
with homogeneous kernel generated by the original relative
relations $g_\nu$, exact origin and weight lifting, and
smooth-point dimensions $n_B(j)+a_j$.
Moreover every quotient-valued point derivation on $U$
whose restriction to $B_0$ lifts to an $O$-valued derivation
of weight $j$, and which is zero on new variables of other
weights, lifts to an $O$-valued point derivation of weight
$j$ on $\mathcal P$.
\end{proposition}
\begin{proof}
First clarify the chart definition.  Put
\[
       V_x=(I/I^2)\otimes_O E.
 \tag{MX.15}
\]
A source \texttt{IsChart} is a homogeneous $O$-algebra map
$\phi:R\to E[[t_1,\ldots,t_n]]$ whose constant coefficient
is $x$ and whose linear-coefficient map identifies $V_x$
with $E^n$.  Indeed its \texttt{inj} clause says precisely
that the kernel on $I$ is
\[
 \operatorname{sat}(I^2)=\{f:\exists\,0\ne d\in O, df\in I^2\},
 \tag{MX.16}
\]
and its \texttt{onto} clause supplies each coordinate
direction with a nonzero coefficient.  Conversely an
isomorphism on $V_x$ supplies those directions after
clearing finitely many denominators.  The variables $t_i$
have the specified weights, and their counts are
\texttt{SmoothPt}.  This is a rational cotangent chart;
we make no integral formal-smoothness assertion.

We prove (MX.13), including the injectivity that is needed
for saturation.  The images of the finitely many formal
variables generate $I$.  A series modulo their $n$th power
is its finite Taylor polynomial, which belongs to
$R^{\mathrm{gr}}$; this proves surjectivity.
For injectivity, express any $f\in I^n\cap R^{\mathrm{gr}}$
as a finite sum
$f=\sum_{|a|=n}z^a h_a$, $h_a\in R$.
Let $J_f$ be the finite set of weights of $f$.
Homogeneity of the presentation kernel makes each weight
component of an element of $R$ well defined.  Taking the
$J_f$-components of this equation replaces $h_a$ by the
finite sum of its components of weights
$j-\operatorname{wt}(z^a)$, $j\in J_f$.
The new coefficients belong to $R^{\mathrm{gr}}$ and give
the same $f$.  Each $z^a$ is a product of $n$ elements
of $I_{\mathrm{gr}}$, proving the contraction equality.

A chart on the homogeneous subalgebra sends
$I_{\mathrm{gr}}^n$ into $(t)^n$.  The isomorphisms (MX.13)
therefore extend it coefficient by coefficient to $R$:
to calculate a coefficient of degree $d$, choose an
approximation modulo $I^{d+1}$.  The answer is independent
of that choice, and addition and multiplication follow
from the corresponding finite-jet operations.
The extension preserves weights and constant coefficients;
the $n=2$ isomorphism proves its tangent injectivity and
surjectivity.  Restricting a chart to $R^{\mathrm{gr}}$
uses the same argument in reverse.  This proves the first
assertion without identifying $R$ with its graded subalgebra.

Write the framed presentation as
\[
 U=O[[z_s,u_t]]/
       (J_B,\,f_\nu),\qquad
 B_0=O[[z_s]]/J_B.
 \tag{MX.17}
\]
Here $J_B$ is extended by the actual old-variable inclusion,
and $f_\nu$ are the relative homogeneous equations from FP8.
Define $\mathcal P=O[[z_s,u_t]]/(J_B)$, with $g_\nu$
the images of $f_\nu$.  This is $B_0[[u_t]]$: coefficientwise
lifts give surjectivity, and finite generation of $J_B$
identifies its coefficientwise extension with its generated
ideal.  The finite power-series ring and its proper quotient
are complete Noetherian local with residue field $k$.
These standard completeness statements are also used in
FP1; the quotient assertion is Stacks, Lemma 10.160.2,
tag 0325.  The origin is evaluation at $u=0$ followed by
$x_B$.  Thus $p$ is surjective, its kernel is exactly
$(g_\nu)$, it preserves origins, and every homogeneous
element of $U$ has a homogeneous preimage in $\mathcal P$.

Extend the given chart of $B_0^{\mathrm{gr}}$ to $B_0$
using (MX.13), and send each new variable to a new formal
variable of the same weight.  This gives a chart of
$\mathcal P$ with dimensions $n_B(j)+a_j$.
The cotangent assertion follows either directly from
coefficients or from
\[
 V_{x_{\mathcal P}}=V_{x_B}\oplus
                              \bigoplus_{t\in\tau} E[u_t].
 \tag{MX.18}
\]

For the derivation assertion let
$D_U:U\to O/(a)$ and $D_B:B_0\to O$ be the given point
derivations, with $D_Ui=D_B\bmod a$.
For new variables of weight $j$ choose any integral lifts
$d_t$ of $D_U(pu_t)$; put $d_t=0$ for all other weights.
There is a unique point derivation $\delta$ on
$B_0[[u_t]]$ restricting to $D_B$ and sending $u_t$ to
$d_t$.  Explicitly, if $F=b_0+\sum_t b_tu_t+$ terms of
degree at least two in $u$, then
\[
       \delta(F)=D_B(b_0)+\sum_t x_B(b_t)d_t.
 \tag{MX.19}
\]
This proves the Leibniz identity even for arbitrary formal
series.  The assigned values and the weight of $D_B$
make $\delta$ homogeneous of weight $j$.
Modulo $(a)$ it equals $D_Up$: two point derivations on a
finite formal-series presentation agreeing on its variables
agree everywhere, since both kill the square of the
augmentation ideal.  This is every field of the cover's
derivation-lifting clause.
\end{proof}
% END MX4

% BEGIN MX5
\begin{proposition}[MX5: independent homogeneous cuts with tangent preservation]
\label{claim:MX5}
Let an augmented weighted $\mathcal P\in\mathcal C_O$
have a source chart, and let $t_1,\ldots,t_h$ be homogeneous
elements of origin zero.  Suppose each $t_i$ has an
$O$-valued point derivation nonzero on $t_i$ and zero on
the other $t_j$.  There is a surjection
\[
        q:\mathcal P\twoheadrightarrow R
 \tag{MX.20}
\]
onto a domain in $\mathcal C_O$, killing every $t_i$,
preserving origins and homogeneous lifting, and carrying
a chart with one variable of weight $\operatorname{wt}(t_i)$
removed for each $i$.  Every $O$-valued point derivation on
$\mathcal P$ that kills all $t_i$ annihilates $\ker q$.
The final chart also restricts to the homogeneous subalgebra
of $R$, exactly as required by \texttt{Model.smooth}.
\end{proposition}
\begin{proof}
Every point derivation kills $I_{\mathcal P}^2$ and,
because $O$ is a domain, its saturation.  Applying the
separating derivations to a relation between the classes
of the $t_i$ shows those classes are $O$-linearly independent
modulo the saturated square.  Clearing denominators makes
them independent in $V_{x_{\mathcal P}}$ as well.
Their chart images therefore have independent linear terms
over $E$, with the specified weights.

Within each weight block perform Gaussian elimination on
these linear terms.  Choose $h$ pivot variables and
invert the corresponding matrix over $E$; no inversion in
$O$ is asserted.  The resulting equations have the form
\[
       x_i+\sum_j c_{ij}y_j+H_i(x,y)=0,
                   \qquad H_i\in(x,y)^2,
 \tag{MX.21}
\]
homogeneously in each original weight.  They can be solved
uniquely in $E[[y]]$ by successive homogeneous degree:
the degree-one part is $-\sum_jc_{ij}y_j$, and the
degree-$n$ part is determined by the already known lower
parts of $-H_i$.  Equivalently iteration
$x^{(n+1)}=-Cy-H(x^{(n)},y)$ converges in the variable-adic
topology.  Every iteration preserves weights.
Substitution into the original chart gives
$\chi:\mathcal P\to E[[y]]$ killing the $t_i$.

First put $Q=\mathcal P/(t_i)$.  Its rational cotangent
is $V_{x_{\mathcal P}}$ modulo the span of the classes of
the $t_i$.  The linear map induced by (MX.21) identifies
that quotient with the remaining coordinate space.
Hence the induced $\chi_Q:Q\to E[[y]]$ is still a source
chart.  Now put
\[
                  R=Q/\ker\chi_Q.
 \tag{MX.22}
\]
This is a domain, since it embeds as a subring of $E[[y]]$.
The kernel lies in the augmentation ideal, by the constant
coefficient identity, so the origin descends and the
quotient is nontrivial.  Proper quotients of complete
Noetherian local rings are complete Noetherian local;
their residue field remains $k$ here.  Thus $R\in\mathcal C_O$.

The entire kernel of $O[[z_s]]\to E[[y]]$ is homogeneous.
Indeed the images of the $z_s$ have zero constant term, so
an output coefficient of degree $n$ sees only finitely
many input monomials of degree at most $n$.  On this finite
sum the weight decomposition commutes with the map.
If its whole image is zero, the image of each weight
component is therefore zero coefficient by coefficient.
Consequently $R$ has the induced weighted presentation,
with no new variable weights, and its weight pieces are
exactly the images of those of $\mathcal P$.

To check the last, essential assertion, let $\delta$ kill
all $t_i$.  The Leibniz rule and $x(t_i)=0$ show that it
kills their entire ideal, so it descends to $Q$.
If $f\in\ker\chi_Q$, the chart property says
$df\in I_Q^2$ for some $0\ne d\in O$.
Then $d\delta(f)=0$, and $\delta(f)=0$ because $O$ is
a domain.  Thus it descends further through (MX.22).
The new chart retains its tangent isomorphism: taking
this last quotient removes only classes already zero in
the rational cotangent.  Finally apply (MX.13) for $n=2$
to restrict it to $R^{\mathrm{gr}}$.
This proves the domain, all exact dimension losses, and
the integral derivation preservation simultaneously.
\end{proof}
% END MX5

\subsection{Normalization and the final rings}

% BEGIN MX6
\begin{proposition}[MX6: independent six-root cuts and exact normalization]
\label{claim:MX6}
In the actual framed ring $U$ define the six positions
\[
 \mathcal P_6=((0,1),(0,3),(2,1),(1,0),(3,0),(1,2)).
 \tag{MX.23}
\]
For $v=\ell$ take the entries in these positions of
$r_U(D.ph)$; for $v\in S_b$ take those of $\alpha_v$.
Denote the resulting $6(b+1)$ elements by $t_{v,h}$.
They have origin zero, have the six weights
$A,H,-B,-A,-H,B$, and admit the separating integral
derivations needed in MX5.
After killing them and, for each $v\in S_b$, the entries
$(0,2),(2,0),(0,0),(1,1)$ of $\alpha_v-I$, every
$\alpha_v$ is exactly $I$.
The six cuts at $\ell$ make $r(D.ph)$ block diagonal in
the source's $\lambda$-blocks.
\end{proposition}
\begin{proof}
We first construct actual tangent directions.
For any $Y\in\mathfrak {sp}_4(O)$ and
$A_v\in\mathfrak {sp}_4(O)$ with $A_\ell=0$, conjugate
the diagonal global representation over $O[\epsilon]$
by $I+\epsilon Y$.  Its derivative is $[Y,r_0(g)]$.
Take $\alpha'_v=I+\epsilon A_v$ and conjugate the
diagonal local tuple by $I+\epsilon(Y-A_v)$.
The local equations hold: the symplectic form is preserved
because $Y^*=-Y$, and characteristic polynomials,
polynomial projectors, their products and inertia identities
are invariant under conjugation.  The diagonal local
equations were proved in MX1.
The universality of $B_0$ gives a map to $O[\epsilon]$
for this local tuple.  Direct multiplication with
$\epsilon^2=0$ gives
\[
 r'(g)(I+\epsilon A_v)
  =(I+\epsilon A_v)\rho'_v(g).
 \tag{MX.24}
\]
The global multiplier and \texttt{UnrT} conditions also
hold: conjugating a matrix equal to $I$ leaves it equal
to $I$.  Thus MX3 gives a map $U\to O[\epsilon]$.
Its constant part is $x_U$, by the two universality
statements, so its $\epsilon$-part is a point derivation
with values
\[
 \delta(r_U(g))=[Y,r_0(g)],\qquad
                         \delta(\alpha_v)=A_v.
 \tag{MX.25}
\]

For the six positions in (MX.23), respectively, take the
symplectic root matrices
\[
 E_{01}-E_{23},\ E_{03},\ E_{21},\
 E_{10}-E_{32},\ E_{30},\ E_{12}.
 \tag{MX.26}
\]
Their six chosen entries form the standard basis.
For a cut at $v\ne\ell$, use $Y=0$, that root matrix
for $A_v$, and zero for the other $A_{v'}$.
For a cut at $\ell$, use that root matrix for $Y$ and
all $A_v=0$.
The derivative of the selected entry of $r_U(D.ph)$ is
the difference of its two diagonal characters at $D.ph$,
and the other five selected entries have derivative zero.
AX6 gives $\bar r(D.ph)=\bar r(D.sg(\ell))$ and the
cross-$\lambda$ residual eigenvalue separation from
\texttt{Good.condD}.  The difference is therefore a unit
of $O$, in particular nonzero.  Formula (MX.25) gives
the asserted independent directions for all $t_{v,h}$.
Their origins and weights follow from the actual entries
and the diagonal base tuple.

We record the matrix calculation giving exact
normalization; it is valid in any local ring $A$.
Let $g^*g=\nu I$ with $\nu$ a unit, all off-diagonal
entries in $\mathfrak m_A$, and the six entries (MX.23)
zero.  The $(0,0)$ and $(1,1)$ multiplier equations show
that $g_{00}g_{33}$ and $g_{11}g_{22}$ are units: each
is $\nu$ plus a product containing an off-diagonal
entry.  Thus $g_{00},g_{11}$ are units.
The $(0,1)$ and $(1,0)$ equations then give
$g_{23}=g_{32}=0$.  These, together with the original
six zeros, are exactly all cross-$\lambda$ entries.
If also $\nu=1$, $g_{02}=g_{20}=0$ and
$g_{00}=g_{11}=1$, the same equations give
$g_{22}=g_{33}=1$; the $(1,3)$ and $(2,0)$ equations
then give $g_{13}=g_{31}=0$.  Hence $g=I$.

Apply this calculation to each frame matrix, whose
reduction is $I$, and to $r_U(D.ph)$, whose reduction
is diagonal.  It proves both asserted normalizations
after every subsequent quotient, with no appeal to a
generic characteristic-zero point.
\end{proof}
% END MX6

% BEGIN MX7
\begin{proposition}[MX7: all rings, equations, dimensions and local components]
\label{claim:MX7}
From the data of MX3 one can construct rings $R,S\in\mathcal C_O$,
with $R$ a domain, a weighted presentation $W$ on $R$,
an origin $x_R$, and a surjection $\pi:R\to S$, having
all the \texttt{Model} fields except those concerning
generation and cocycle derivations.
Their dimensions are explicitly
\[
 r=1+b+a_{-P},\qquad r'=2+b+a_P,\qquad
                         m=2+2b+a_0.
 \tag{MX.27}
\]
The kernel of $\pi$ is generated by exactly $r+1$
equations $F$ of weight $-P$, $r'$ equations $E$ of
weight $P$, and $m$ equations $Z$ of weight zero and
origin zero.
\end{proposition}
\begin{proof}
Use the cover $p:\mathcal P\to U$ of MX4 and choose
homogeneous lifts $\widetilde t_{v,h}$ of the cuts in
MX6.  Their origins are zero.  Pulling the separating
derivations back along $p$ gives the hypotheses of MX5.
Apply it to obtain $q:\mathcal P\twoheadrightarrow R$,
with all its domain, origin, weight and tangent properties.

In $U$ also form the following homogeneous elements:
the four entries of $\alpha_v-I$ at positions
$(0,2),(2,0),(0,0),(1,1)$ for each $v\in S_b$,
and the two entries $(0,2),(2,0)$ of
\[
                 h_U=r_U(g_0)+\vartheta(g_0)r_U(g_0^{-1}).
 \tag{MX.28}
\]
Call this list $s_\eta$.  All its origins are zero;
its weights are $P,-P,0,0$ per $v$ and $P,-P$ for
the last two.  Choose homogeneous lifts
$\widetilde s_\eta\in\mathcal P$ and set
\[
 \begin{gathered}
 J_S=(q(g_\nu),q(\widetilde s_\eta))\subset R,
 \qquad S=R/J_S,\qquad\pi:R\twoheadrightarrow S,\\
 \sigma:U\twoheadrightarrow S,
 \qquad \sigma p=\pi q,
 \qquad i_R=q|_{B_0}.
 \end{gathered}
 \tag{MX.29}
\]
The map $\sigma$ exists because $\ker p=(g_\nu)$;
it is surjective because $p,q,\pi$ are.
Every generator of $J_S$ has origin zero, so this ideal
is proper and $S\in\mathcal C_O$.
The same origin descends to $S$ from both sides of
the square.  In particular
\[
 \pi f=\sigma u\quad\Longrightarrow\quad
                            x_R(f)=x_U(u).
 \tag{MX.30}
\]
Every homogeneous $u\in U_j$ has a lift $f\in R_j$
with $\pi f=\sigma u$: lift first through $p$, then
apply $q$.  All presentation-variable weights in $R$
already occurred in $\mathcal P$, hence in $U$ or
$B_0$, and therefore belong to $\mathcal W$.

The smooth-point dimensions before the cuts are
$d_\ell(j)+bd_p(j)+a_j$.
Each of the six nonzero weights $\pm A,\pm B,\pm H$
loses $b+1$ dimensions.  At $A,B,H$ the result is
$2+b-(b+1)=1$, and at their negatives it is zero.
No cut has weight $0,P,-P$, giving respectively
$m,r',r$ in (MX.27).  The resulting dimension function
is exactly \texttt{tDims r r' m}; in particular $r\ge1$.
MX5 already gives this chart on the homogeneous subalgebra
of $R$, with its restricted origin, as the record demands.

The relative equations in (MX.12) occur only at
$0,P,-P$.  The additional list $s_\eta$ contributes
$b+1$ equations at each of $P,-P$ and $2b$ at $0$.
Thus their total counts are
\[
 \begin{array}{c|c|c}
 \text{weight}&\text{number of equations}&\text{dimension}\\ \hline
 -P&a_{-P}+1+b+1=r+1&r\\
 P&a_P+1+b+1=r'&r'\\
 0&a_0+2+2b=m&m.
 \end{array}
 \tag{MX.31}
\]
Enumerate these complete lists, including any zero
equations, as $F,E,Z$.  They generate exactly $J_S$,
so the prescribed kernel equality holds.  Their weights
and the origin statement for $Z$ have already been proved.

Define $r_S=\sigma_*r_U$.  It is a continuous lift of
the original $\bar r$.  Multiplier and \texttt{UnrT}
pass through $\sigma$.  MX6 gives $\sigma(\alpha_v)=I$
and the block condition at $D.ph$; (MX.28) gives the
other two normalization equations.  Hence
\[
      (r_S)|_{G_v}=\pi_* (i_R)_*\rho_v.
 \tag{MX.32}
\]
For each local condition, use this very domain $R$ as
the \texttt{OnComp} witness.  Its local lift
$(i_R)_*\rho_v$ satisfies the original polynomial
equations, and $x_Ri_R=x_B$ maps it to $r_0|_{G_v}$.
The map to $S$ is exactly $\pi$, so (MX.32) proves
both component clauses of \texttt{D.IsPt}.  There is
no assumption that a solution automatically lies on the
desired component.
Finally homogeneous entry lifting follows from the
lifting statement before (MX.30), and for every preimage
of any matrix entry (MX.30) gives its prescribed value
at the origin.  These are \texttt{hom} and \texttt{at0},
with their universal preimage quantifiers.
\end{proof}
% END MX7

% BEGIN MX8
\begin{proposition}[MX8: the actual representation generates the quotient]
\label{claim:MX8}
The representation $r_S$ of MX7 satisfies the literal
\texttt{Lift.Generates}: for every $A\in\mathcal C_O$,
two $O$-algebra maps $f,f':S\to A$ with
$f_*r_S=f'_*r_S$ are equal.
\end{proposition}
\begin{proof}
By (MX.32), the maps
$f\sigma i,f'\sigma i:B_0\to A$ give the same local
tuple, namely the restrictions of the equal global
representations.  The polynomial local equations are
preserved, so the uniqueness clause of the actual local
base gives $f\sigma i=f'\sigma i$.
Both maps $f\sigma,f'\sigma:U\to A$ now give the
same base map, global representation and frame tuple:
every frame matrix maps to $I$ by MX6.
The uniqueness clause of the actual framed universal
ring gives $f\sigma=f'\sigma$.  Surjectivity of
$\sigma$ proves $f=f'$.  This verifies the source's
test against all rings in $\mathcal C_O$; an informal
statement about a dense algebra of entries is unnecessary.
\end{proof}
% END MX8

\subsection{The integral Selmer-cocycle derivation}

% BEGIN MX9
\begin{proposition}[MX9: every original Sha cocycle gives the required derivation]
\label{claim:MX9}
For every nonzero $a\in O$ and every literal
$z\in\texttt{D.ShaCoc}(a)$, the rings and representation
of MX7 admit an $O$-valued point derivation $\delta$ of
weight $-P$ such that
\[
 \begin{gathered}
 \pi f=0\ \Longrightarrow\ \delta(f)\in(a),\\
 \pi f=r_S(g)_{20}\ \Longrightarrow\
       \delta(f)\bmod a=\vartheta(g)z(g),\\
 \pi f=r_S(g)_{31}\ \Longrightarrow\
       \delta(f)\bmod a=\Psi(g)z(g).
 \end{gathered}
 \tag{MX.33}
\]
All local coboundaries and the normalization $z(g_0)=0$
are those in the original source definition.
\end{proposition}
\begin{proof}
Put $Q=O/(a)$ and $\gamma=\Psi^{-1}$.
The hypotheses say that $z:G\to Q$ is locally constant,
$z(gh)=z(g)+\gamma(g)z(h)$, satisfies
\texttt{UnrT} with value zero, vanishes at $g_0$, and
on each original $G_v$ is a coboundary.
Choose an integral lift $m_v\in O$ of its specified
local parameter, so
\[
             z(g)=\gamma(g)m_v-m_v\pmod a
                                      \quad(g\in G_v).
 \tag{MX.34}
\]
No division by $\gamma(g)-1$ is used.

Work first in the square-zero ring $O\oplus Q$.
The matrix-valued function
\[
 r_z(g)=r_0(g)+\epsilon M r_0(g)z(g),
                    \qquad M=E_{20}+E_{31},
 \tag{MX.35}
\]
is a continuous representation of the same multiplier.
Indeed
$r_0(g)M r_0(g)^{-1}=\gamma(g)M$,
$M^*=-M$, and the cocycle identity proves multiplication
when $\epsilon^2=0$.  Adic continuity follows from that
of $r_0$ and local constancy of $z$.
For residual-trivial inertia away from $T$, both
$r_0(g)=I$ and $z(g)=0$ by MX1 and \texttt{ShaCoc}.
Thus (MX.35) satisfies the actual global unramified
condition as well.

For each $v$ conjugate the original diagonal local
representation over $O[\epsilon]$ by
$I-\epsilon m_vM$.  This satisfies the local equations
by MX6's conjugation argument and has derivative
\[
 [-m_vM,r_0(g)]
       =(\gamma(g)m_v-m_v)Mr_0(g).
 \tag{MX.36}
\]
The local universality gives a map $y_B:B_0\to O[\epsilon]$
for the whole tuple.  Its constant part is $x_B$ by
uniqueness.  Its derivative $D_B$ has weight $-P$.
We justify this last statement on every homogeneous
element, not only on representation entries.

The local entries determine maps from $B_0$ to
$O[\epsilon]$, by the same local uniqueness.
For a unit character $s:\mathbb Z^2\to O^\times$,
the homogeneous presentation defines a torus automorphism
of $B_0$, multiplying weight-$j$ elements by $s(j)$.
Precompose $y_B$ with it and multiply the $\epsilon$
coordinate by $s(-P)^{-1}$.
This map agrees with $y_B$ on every local entry:
the constant part is zero at nonzero weight, and the
derivative in (MX.36) is supported only at weight $-P$.
Thus the two maps agree on all of $B_0$.
For $x\in(B_0)_j$ this gives
\[
           (s(j)/s(-P)-1)D_B(x)=0.
 \tag{MX.37}
\]
If $j\ne-P$, choose $s$ using powers of the
infinite-order unit $1+\ell\in O^\times$ in a coordinate
where $j$ and $-P$ differ.  The parenthesis is nonzero;
the domain property of $O$ forces $D_B(x)=0$.
This establishes the precise \texttt{HasWt} clause.

Reduce the second coordinate of $y_B$ modulo $a$.
Equations (MX.34)--(MX.36) show that its local tuple
is exactly the restriction of $r_z$.  Therefore
$(y_B\bmod a,r_z,(I)_v)$ satisfies (MX.11).
The actual framed universality gives
$y_U:U\to O\oplus Q$.
Its constant part is $x_U$, again by uniqueness, so
its second coordinate is a point derivation $D_U$ with
\[
 \begin{split}
 D_Ui&=D_B\bmod a,\\
 D_U(r_U(g)_{ij})&=M_{ij}\psi_j(g)z(g),\\
 D_U((\alpha_v-I)_{ij})&=0,
 \end{split}
 \tag{MX.38}
\]
where $(\psi_0,\psi_1,\psi_2,\psi_3)
=(\vartheta,\Psi,\alpha,1)$.
Every new presentation variable of weight other than
$-P$ has derivative zero: by MX3 it is a centered
representation or frame entry, and (MX.38) proves the
claim with its original weight.

Apply the integral lifting assertion of MX4 to obtain
an $O$-valued point derivation $\widetilde\delta$ on
$\mathcal P$ of weight $-P$ reducing to $D_Up$.
All cuts $\widetilde t_{v,h}$ have one of the six other
nonzero weights.  Thus $\widetilde\delta$ kills them,
and MX5 makes it descend to an $O$-valued point
derivation $\delta$ on the \emph{actual domain} $R$,
still of weight $-P$.

It remains to check the later quotient $S$, where an
integral descent is not claimed.  The old relations have
$D_U$-value zero.  The frame equations in $s_\eta$
have value zero by (MX.38).  The cocycle identity gives
$z(1)=0$, and then $z(g_0)=0$ gives $z(g_0^{-1})=0$.
Both \texttt{hEl} equations in (MX.28) consequently
have $D_U$-value zero as well.
Thus $\delta\bmod a$ kills every generator of $J_S$.
All these generators have origin zero, so the Leibniz
rule makes it kill their entire ideal.

More generally, for every $f\in R,u\in U$ with
$\pi f=\sigma u$, choose lifts $f=q(h)$ and $u=p(h')$
in $\mathcal P$.  Their difference $q(h-h')$ belongs
to $J_S$.  The preceding vanishing and the reduction
$\widetilde\delta\bmod a=D_Up$ give
\[
                  \delta(f)\bmod a=D_U(u).
 \tag{MX.39}
\]
This proves the first clause of (MX.33) by taking $u=0$.
For the other two take the actual $(2,0)$ and $(3,1)$
entries of $r_U(g)$ and use (MX.38): $M_{20}=M_{31}=1$,
$\psi_0=\vartheta$, and $\psi_1=\Psi$.
The equalities hold for every preimage $f$, as required.
This finishes the entire \texttt{Model.der} field.
\end{proof}
% END MX9

% BEGIN MX10
\begin{theorem}[MX10: the literal actual Model constructor]
\label{claim:MX10}
On every original $\Omega,\ell,c,D$ the following exact
implication holds, with the universal \texttt{LemmaP}
already supplied by FP8:
\[
 \begin{gathered}
 \texttt{Good D}\ \wedge\ \texttt{PhiCond D.Gam}\
 \wedge\ \texttt{D.Nums}\ \wedge\
                         \texttt{Nonempty(LocBase D)}\\
             \Longrightarrow\quad\texttt{Nonempty(Model D)}.
 \end{gathered}
 \tag{MX.40}
\]
Consequently DX2, DX4 and DX8, in their exact scopes on
the same good datum, give
\[
       \texttt{Good D}\quad\Longrightarrow\quad
                              \texttt{Nonempty(Model D)}.
 \tag{MX.41}
\]
This conclusion includes every field of the record,
including its domain, actual local components, generation
and integral cocycle derivations.
\end{theorem}
\begin{proof}
Choose the supplied local base once.  MX1--MX3 give its
actual framed universal deformation ring with its origin,
homogeneous variables, relative equations and exact counts.
MX4 constructs the relative smooth cover.  MX5--MX7
construct the domain $R$, the quotient $S$, the presentation,
origin, integers $(r,r',m)$, all three equation lists,
and the actual representation $r_S$ with its \texttt{IsPt}
witnesses.  MX8 supplies generation and MX9 supplies
the complete derivation assertion.
The following table identifies every record field.
\[
 \begin{array}{c|l}
 \text{source fields}&\text{constructed supplier}\\ \hline
 R,\ \texttt{IsCO},\ \texttt{IsDomain},\ W,\ \texttt{org}
       &\text{MX4--MX7}\!\\
 \texttt{wt\_mem},\ r,r',m,\ \texttt{one\_le},\ \texttt{smooth}
       &\text{MX7, (MX.27)}\\
 F,E,Z,\ \texttt{hF},\texttt{hE},\texttt{hZ}
       &\text{MX7, (MX.31)}\\
 S,\ \texttt{IsCO},\ \pi,\ \texttt{surj},\ \texttt{ker}
       &\text{MX7, (MX.29)}\\
 \texttt{ru},\ \texttt{pt},\ \texttt{hom},\ \texttt{at0}
       &\text{MX7, (MX.30)--(MX.32)}\\
 \texttt{gen}&\text{MX8}\\
 \texttt{der}&\text{MX9, (MX.33).}
 \end{array}
 \tag{MX.42}
\]
All instances and coefficient maps belong to the chosen
rings; every local component is witnessed by the same
constructed domain $R$.  This is exactly the source
\texttt{exModel} conclusion and proves (MX.40).
For (MX.41), substitute DX2 for \texttt{PhiCond}, DX4
for \texttt{Nums}, and the actual local base of DX8.
No weakening of those predicates, change of original
coefficient data, or converse to a conditional theorem
is used in this substitution.
\end{proof}
% END MX10

\paragraph{Scope of the completed construction.}
The new result supplies an actual \texttt{Model} on the
original good datum.  Its construction no longer leaves
\texttt{Fr}, \texttt{Cover}, a generic smooth quotient,
normalization, generation or a cocycle-derivation package
as an unexplained supplier.  The local equations enter
through the actual \texttt{LocBase}; the arithmetic
cohomology enters through the precise \texttt{Nums} and
\texttt{PhiCond} interfaces, whose installations are DX2,
DX4 and DX8.  The first-weight-zero geometry and the
representation-point/orbit/trace comparison are addressed
separately by PX.  None of these constructions by itself
asserts automorphy, image density, rigidity, the finite
trace catalogue, an unconditional count statement, or the
complete RH/Dirichlet-GRH reconstruction.
